\frfilename{mt2a2.tex}
\versiondate{30.11.09}
\copyrightdate{2000}

\def\chaptername{Lampiran}
\def\sectionname{Topologi ruang Euklides}

\newsection{2A2}

Dalam lampiran Jilid 1\cmmnt{ (\S1A2)} saya membahas himpunan terbuka dan
tertutup dalam $\BbbR^r$; tujuan utama di sana adalah menopang gagasan
`himpunan Borel', yang sangat penting dalam teori ukuran Lebesgue, tetapi
tentu saja himpunan-himpunan itu juga mendasar bagi kajian fungsi kontinu,
dan bahkan bagi semua aspek analisis real.
Di sini saya memberikan pengantar yang sangat singkat mengenai fakta-fakta
elementer lebih lanjut tentang himpunan tertutup dan kompak serta fungsi
kontinu yang kita perlukan untuk jilid ini. Sebagian besar materi ini dapat
diturunkan dari perumuman dalam \S2A3, tetapi saya tetap menguraikan garis
besar buktinya, sebab untuk sebagian besar jilid ini (kebanyakan
pengecualiannya berada dalam Bab 24) ruang Euklides sudah memadai bagi
keperluan kita.

\leader{2A2A}{Penutupan: Definisi} Untuk setiap $r\ge 1$ dan setiap
$A\subseteq\BbbR^r$, {\bf penutupan} $A$, $\overline{A}$, adalah irisan
semua subhimpunan tertutup dari $\BbbR^r$ yang memuat $A$.
Ini adalah\cmmnt{ himpunan tertutup (karena merupakan irisan suatu keluarga
tak kosong dari himpunan-himpunan tertutup, lihat 1A2Fd), sehingga merupakan}
himpunan tertutup terkecil yang memuat $A$. \cmmnt{Khususnya, }$A$ tertutup
jika dan hanya jika $\overline{A}=A$.

\leader{2A2B}{Lema} Misalkan $A\subseteq\BbbR^r$ sebarang himpunan. Maka,
untuk $x\in\BbbR^r$, pernyataan-pernyataan berikut ekuivalen:

(i) $x\in\overline{A}$\cmmnt{, penutupan $A$};

(ii) $B(x,\delta)\cap A\ne\emptyset$ untuk setiap $\delta>0$, dengan
$B(x,\delta)=\{y:\|y-x\|\le\delta\}$;

(iii) terdapat suatu barisan $\sequencen{x_n}$ dalam $A$ sedemikian sehingga
$\lim_{n\to\infty}\|x_n-x\|=0$.

\proof{{\bf (a)(i)$\Rightarrow$(ii)} Andaikan $x\in\overline{A}$ dan
$\delta>0$. Maka $U(x,\delta)=\{y:\|y-x\|<\delta\}$ adalah himpunan terbuka
(1A2D), sehingga $F=\BbbR^r\setminus U(x,\delta)$ tertutup, sedangkan
$x\notin F$. Sekarang

\Centerline{$x\in\overline{A}\setminus F
\Longrightarrow\overline{A}\not\subseteq F
\Longrightarrow A\not\subseteq F
\Longrightarrow A\cap U(x,\delta)\ne\emptyset
\Longrightarrow A\cap B(x,\delta)\ne\emptyset$.}

\noindent Karena $\delta$ sebarang, (ii) benar.

\medskip

{\bf (b)(ii)$\Rightarrow$(iii)} Jika (ii) benar, maka untuk setiap $n\in\Bbb
N$ kita dapat menemukan $x_n\in A$ sedemikian sehingga
$\|x_n-x\|\le 2^{-n}$, dan sekarang $\lim_{n\to\infty}\|x_n-x\|=0$.

\medskip

{\bf (c)(iii)$\Rightarrow$(i)} Andaikan (iii). \Quer\ Andaikan, jika
mungkin, bahwa $x\notin\overline{A}$. Maka $x$ termasuk dalam himpunan terbuka
$\BbbR^r\setminus\overline{A}$ dan terdapat $\delta>0$ sedemikian sehingga
$U(x,\delta)\subseteq\BbbR^r\setminus\overline{A}$. Tetapi sekarang terdapat
$n$ sedemikian sehingga $\|x_n-x\|<\delta$, dan dalam hal ini
$x_n\in U(x,\delta)\cap A\subseteq U(x,\delta)\cap\overline{A}$. \Bang
}%end of proof of 2A2B

\leader{2A2C}{Fungsi kontinu (a)} \cmmnt{Saya mulai dengan suatu
karakterisasi fungsi kontinu dalam istilah himpunan terbuka.} Jika $r$,
$s\ge 1$, $D\subseteq\BbbR^r$ dan $\phi:D\to\BbbR^s$ adalah suatu fungsi,
kita katakan bahwa $\phi$ {\bf kontinu} jika untuk setiap $x\in D$ dan
$\epsilon>0$ terdapat $\delta>0$ sedemikian sehingga
$\|\phi(y)-\phi(x)\|\le\epsilon$ kapan pun $y\in D$ dan
$\|y-x\|\le\delta$. \cmmnt{Sekarang }$\phi$ kontinu jika dan hanya jika
untuk setiap himpunan terbuka $G\subseteq\BbbR^s$ terdapat suatu himpunan
terbuka $H\subseteq\BbbR^r$ sedemikian sehingga
$\phi^{-1}[G]=D\cap H$.

\prooflet{\Prf\ {\bf (i)} Andaikan $\phi$ kontinu dan
$G\subseteq\BbbR^s$ terbuka. Tetapkan

\Centerline{$H=\bigcup\{U:U\subseteq\BbbR^r\text{ terbuka},\,\phi[U\cap
D]\subseteq G\}$.}

\noindent Maka $H$ merupakan gabungan himpunan-himpunan terbuka, sehingga
terbuka (1A2Bd), dan $H\cap D\subseteq\phi^{-1}[G]$. Jika
$x\in\phi^{-1}[G]$, maka $\phi(x)\in G$, sehingga terdapat $\epsilon>0$
sedemikian sehingga $U(\phi(x),\epsilon)\subseteq G$; sekarang terdapat
$\delta>0$ sedemikian sehingga
$\|\phi(y)-\phi(x)\|\le\bover12\epsilon$ kapan pun $y\in D$ dan
$\|y-x\|\le\delta$, sehingga

\Centerline{$\phi[U(x,\delta)\cap D]
\subseteq U(\phi(x),\epsilon)\subseteq G$}

\noindent dan

\Centerline{$x\in U(x,\delta)\subseteq H$.}

\noindent Karena $x$ sebarang, $\phi^{-1}[G]=H\cap D$. Karena $G$
sebarang, $\phi$ memenuhi syarat tersebut.

\quad{\bf (ii)} Sekarang andaikan bahwa $\phi$ memenuhi syarat tersebut.
Ambil $x\in D$ dan $\epsilon>0$. Maka $U(\phi(x),\epsilon)$ terbuka,
sehingga terdapat himpunan terbuka $H\subseteq\BbbR^r$ sedemikian sehingga
$H\cap D=\phi^{-1}[U(\phi(x),\epsilon)]$; kita melihat bahwa $x\in H$,
sehingga terdapat $\delta>0$ sedemikian sehingga
$U(x,\delta)\subseteq H$; sekarang jika $y\in D$ dan
$\|y-x\|\le\bover12\delta$, maka $y\in D\cap H$,
$\phi(y)\in U(\phi(x),\epsilon)$ dan
$\|\phi(y)-\phi(x)\|\le\epsilon$. Karena $x$ dan $\epsilon$ sebarang,
$\phi$ kontinu. \Qed}

\header{2A2Cb}{\bf (b)} Dengan menggunakan definisi kontinuitas
$\epsilon$-$\delta$, mudah dilihat bahwa suatu fungsi $\phi$ dari
subhimpunan $D$ dari $\BbbR^r$ ke $\BbbR^s$ kontinu jika dan hanya jika semua
komponennya $\phi_i$ kontinu, dengan menuliskan
$\phi(x)=(\phi_1(x),\ldots,\phi_s(x))$ untuk $x\in D$.
\prooflet{\Prf\
{\bf (i)} Jika $\phi$ kontinu, $i\le s$, $x\in D$ dan $\epsilon>0$, maka
terdapat $\delta>0$ sedemikian sehingga

\Centerline{$|\phi_i(y)-\phi_i(x)|\le\|\phi(y)-\phi(x)\|\le\epsilon$}

\noindent kapan pun $y\in D$ dan $\|y-x\|\le\delta$. {\bf (ii)} Jika
setiap $\phi_i$ kontinu, $x\in D$ dan $\epsilon>0$, maka terdapat
$\delta_i>0$ sedemikian sehingga
$|\phi_i(y)-\phi_i(x)|\le\epsilon/\sqrt{s}$ kapan pun $y\in D$ dan
$\|y-x\|\le\delta_i$; dengan menetapkan
$\delta=\min_{1\le i\le s}\delta_i>0$, kita memperoleh
$\|\phi(y)-\phi(x)\|\le\epsilon$ kapan pun $y\in D$ dan
$\|y-x\|\le\delta$. \Qed}

\spheader 2A2Cc\dvAnew{2015}\cmmnt{ Pada satu atau dua tempat, kita akan
menjumpai penguatan berikut dari gagasan `fungsi kontinu'.}
Jika $r$, $s\ge 1$, $D\subseteq\BbbR^r$ dan
$\phi:D\to\BbbR^s$ adalah suatu fungsi, kita katakan bahwa $\phi$
{\bf kontinu seragam} jika untuk setiap $\epsilon>0$ terdapat $\delta>0$
sedemikian sehingga $\|\phi(y)-\phi(x)\|\le\epsilon$ kapan pun $x$,
$y\in D$ dan $\|y-x\|\le\delta$. Suatu fungsi kontinu seragam
\cmmnt{tentu saja }kontinu.

\leader{2A2D}{Kekompakan dalam $\BbbR^r$: Definisi} Suatu subhimpunan $F$
dari $\BbbR^r$ disebut {\bf kompak} jika, setiap kali $\Cal G$ merupakan
suatu keluarga himpunan terbuka yang menutupi $F$, terdapat subhimpunan
berhingga $\Cal G_0$ dari $\Cal G$ yang masih menutupi $F$.

\leader{2A2E}{Sifat-sifat elementer himpunan kompak} Ambil sebarang
$r\ge 1$, serta subhimpunan-subhimpunan $D$, $F$, $G$ dan $K$ dari
$\BbbR^r$.

\header{2A2Ea}{\bf (a)} Jika $K$ kompak dan $F$ tertutup, maka $K\cap F$
kompak. \prooflet{\Prf\ Misalkan $\Cal G$ suatu selimut terbuka dari
$F\cap K$. Maka $\Cal G\cup\{\BbbR^r\setminus F\}$ merupakan selimut
terbuka dari $K$, sehingga mempunyai subselimut berhingga, sebutlah
$\Cal G_0$. Sekarang $\Cal G_0\setminus\{\BbbR^r\setminus F\}$ adalah
subhimpunan berhingga dari $\Cal G$ yang menutupi $K\cap F$. Karena
$\Cal G$ sebarang, $K\cap F$ kompak. \Qed}

\header{2A2Eb}{\bf (b)} Jika $s\ge 1$, $\phi:D\to\BbbR^s$ adalah fungsi
kontinu, $K$ kompak dan $K\subseteq D$, maka $\phi[K]$ kompak.
\prooflet{\Prf\ Misalkan $\Cal V$ suatu selimut terbuka dari $\phi[K]$.
Misalkan $\Cal H$ adalah

\Centerline{$\{H:H\subseteq\BbbR^r\text{ terbuka},\,\exists\enskip
V\in\Cal V,\,\phi^{-1}[V]=D\cap H\}$.}

\noindent Jika $x\in K$, maka $\phi(x)\in\phi[K]$, sehingga terdapat
$V\in\Cal V$ sedemikian sehingga $\phi(x)\in V$; sekarang terdapat
$H\in\Cal H$ sedemikian sehingga $D\cap H=\phi^{-1}[V]$ memuat $x$ (2A2Ca);
karena $x$ sebarang, $K\subseteq\bigcup\Cal H$. Misalkan $\Cal H_0$
subhimpunan berhingga dari $\Cal H$ yang menutupi $K$. Untuk setiap
$H\in\Cal H_0$, misalkan $V_H\in\Cal V$ sedemikian sehingga
$\phi^{-1}[V_H]=D\cap H$; maka $\{V_H:H\in\Cal H_0\}$ adalah subhimpunan
berhingga dari $\Cal V$ yang menutupi $\phi[K]$. Karena $\Cal V$ sebarang,
$\phi[K]$ kompak. \Qed}

\spheader 2A2Ec Jika $K$ kompak, maka himpunan itu tertutup.
\prooflet{\Prf\ Tuliskan $H=\BbbR^r\setminus K$. Ambil sebarang $x\in H$.
Maka $G_n=\BbbR^r\setminus B(x,2^{-n})$ terbuka untuk setiap $n\in\Bbb N$
(1A2G). Juga

\Centerline{$\bigcup_{n\in\Bbb N}G_n=\{y:y\in\BbbR^r,\,\|y-x\|>0\}
=\Bbb R^r\setminus\{x\}\supseteq K$.}

\noindent Jadi terdapat suatu himpunan berhingga
$\Cal G_0\subseteq\{G_n:n\in\Bbb N\}$ yang menutupi $K$.
Harus terdapat $n$ sedemikian sehingga
$\Cal G_0\subseteq\{G_i:i\le n\}$, sehingga

\Centerline{$K\subseteq\bigcup\Cal G_0\subseteq\bigcup_{i\le n}G_i=G_n$,}

\noindent dan $B(x,2^{-n})\subseteq H$. Karena $x$ sebarang, $H$ terbuka
dan $K$ tertutup. \Qed}

\spheader 2A2Ed Jika $K$ kompak dan $G$ terbuka serta $K\subseteq G$, maka
terdapat $\delta>0$ sedemikian sehingga
$K+B(\tbf{0},\delta)\subseteq G$.
\prooflet{\Prf\ Jika $K=\emptyset$, pernyataan ini trivial, sebab dalam
hal ini

\Centerline{$K+B(\tbf{0},1)=\{x+y:x\in K,\,y\in B(\tbf{0},1)\}=\emptyset$.}

\noindent Jika tidak, tetapkan

\Centerline{$\Cal G
=\{U(x,\delta):x\in \Bbb R^r,\,\delta>0,\,U(x,2\delta)\subseteq G\}$.}

\noindent Maka $\Cal G$ adalah suatu keluarga himpunan terbuka dan
$\bigcup\Cal G=G$ (karena $G$ terbuka), sehingga $\Cal G$ merupakan
selimut terbuka dari $K$ dan mempunyai subselimut berhingga $\Cal G_0$.
Nyatakan $\Cal G_0$ sebagai
$\{U(x_0,\delta_0),\ldots,U(x_n,\delta_n)\}$, dengan
$U(x_i,2\delta_i)\subseteq G$ untuk setiap $i$. Tetapkan
$\delta=\min_{i\le n}\delta_i>0$. Jika $x\in K$ dan
$y\in B(\tbf{0},\delta)$, maka terdapat $i\le n$ sedemikian sehingga
$x\in U(x_i,\delta_i)$; sekarang

\Centerline{$\|(x+y)-x_i\|\le\|x-x_i\|+\|y\|<\delta_i+\delta
\le 2\delta_i$,}

\noindent sehingga $x+y\in U(x_i,2\delta_i)\subseteq G$. Karena $x$ dan
$y$ sebarang, $K+B(\tbf{0},\delta)\subseteq G$. \Qed}

\cmmnt{
\medskip

\noindent{\bf Catatan} Hasil ini merupakan bentuk sederhana dari
{\bf lema selimut Lebesgue}.
}

\leader{2A2F}{}\cmmnt{ Nilai gagasan `kekompakan' sangat bertambah oleh
kenyataan bahwa terdapat karakterisasi efektif bagi subhimpunan-subhimpunan
kompak dari $\BbbR^r$.

\medskip

\noindent}{\bf Teorema} Untuk setiap $r\ge 1$, suatu subhimpunan $K$ dari
$\BbbR^r$ kompak jika dan hanya jika tertutup dan terbatas.

\proof{{\bf (a)} Andaikan $K$ kompak. Menurut 2A2Ec, himpunan itu tertutup.
Untuk melihat bahwa himpunan itu terbatas, pertimbangkan
$\Cal G=\{U(\tbf{0},n):n\in\Bbb N\}$. $\Cal G$ seluruhnya terdiri atas
himpunan terbuka, dan $\bigcup\Cal G=\Bbb R^r\supseteq K$, sehingga
terdapat suatu $\Cal G_0\subseteq\Cal G$ berhingga yang menutupi $K$.
Harus terdapat $n$ sedemikian sehingga
$\Cal G_0\subseteq\{U(\tbf{0},i):i\le n\}$, sehingga

\Centerline{$K\subseteq\bigcup\Cal G_0\subseteq\bigcup_{i\le
n}U(\tbf{0},i)=U(\tbf{0},n)$,}

\noindent dan $K$ terbatas.

\medskip

{\bf (b)} Jadi tersisa konversnya; saya harus menunjukkan bahwa suatu
himpunan tertutup terbatas bersifat kompak. Bagian utama argumen ini adalah
bukti dengan induksi pada $r$ bahwa interval tertutup
$[-\tbf{n},\tbf{n}]$ kompak untuk setiap $n\in\Bbb N$, dengan menuliskan
$\tbf{n}=(n,\ldots,n)\in\Bbb R^r$.

\medskip

\quad{\bf (i)} Jika $r=1$ dan $n\in\Bbb N$, serta $\Cal G$ adalah suatu
keluarga himpunan terbuka dalam $\Bbb R$ yang menutupi $[-n,n]$, tetapkan

\Centerline{$A=\{x:x\in[-n,n],\text{ terdapat suatu }\Cal
G_0\subseteq\Cal G\text{ berhingga sedemikian sehingga }
[-n,x]\subseteq\bigcup\Cal G_0\}$.}

\noindent Maka $-n\in A$, sebab jika $-n\in G\in\Cal G$, maka
$[-n,-n]\subseteq\bigcup\{G\}$, dan $A$ terbatas di atas oleh $n$, sehingga
$c=\sup A$ ada dan termasuk dalam $[-n,n]$.

Selanjutnya, $c\in[-n,n]\subseteq\bigcup\Cal G$, sehingga terdapat
$G\in\Cal G$ yang memuat $c$. Misalkan $\delta>0$ sedemikian sehingga
$U(c,\delta)\subseteq G$.
Terdapat $x\in A$ sedemikian sehingga $x\ge c-\delta$. Misalkan $\Cal G_0$
suatu subhimpunan berhingga dari $\Cal G$ yang menutupi $[-n,x]$. Maka
$\Cal G_1=\Cal G_0\cup\{G\}$ adalah subhimpunan berhingga dari $\Cal G$
yang menutupi $[-n,c+{1\over2}\delta]$.
Tetapi $c+{1\over2}\delta\notin A$, sehingga
$c+{1\over2}\delta>n$, dan $\Cal G_1$ adalah subhimpunan berhingga dari
$\Cal G$ yang menutupi $[-n,n]$. Karena $\Cal G$ sebarang, $[-n,n]$
kompak dan induksi pun dimulai.

\medskip

\quad{\bf (ii)} Untuk langkah induksi menuju $r+1$, pandang interval
tertutup $F=[-\tbf{n},\tbf{n}]$, diambil dalam $\BbbR^{r+1}$, sebagai
produk interval tertutup $E=[-\tbf{n},\tbf{n}]$, diambil dalam
$\BbbR^r$, dengan interval tertutup $[-n,n]\subseteq\Bbb R$; menurut
hipotesis induksi, baik $E$ maupun $[-n,n]$ kompak. Misalkan $\Cal G$ suatu
keluarga subhimpunan terbuka dari $\BbbR^{r+1}$ yang menutupi $F$. Tuliskan
$\Cal H$ untuk keluarga subhimpunan terbuka $H$ dari $\BbbR^r$ sedemikian
sehingga $H\times[-n,n]$ ditutupi oleh suatu subkeluarga berhingga dari
$\Cal G$. Maka $E\subseteq\bigcup\Cal H$. \Prf\ Ambil $x\in E$. Tetapkan

\Centerline{$\Cal U_x=\{U:U\subseteq\Bbb R\text{ terbuka},\,\exists\enskip
G\in\Cal G,\text{ terbuka }H\subseteq\BbbR^r,\,
x\in H\text{ dan }H\times U\subseteq G\}$.}

\noindent Maka $\Cal U_x$ adalah suatu keluarga subhimpunan terbuka dari
$\Bbb R$. Jika $\xi\in[-n,n]$, terdapat $G\in\Cal G$ yang memuat
$(x,\xi)$; terdapat $\delta>0$ sedemikian sehingga
$U((x,\xi),\delta)\subseteq G$; sekarang $U(x,\bover12\delta)$ dan
$U(\xi,\bover12\delta)$ masing-masing merupakan himpunan terbuka dalam
$\BbbR^r$, $\Bbb R$, serta

\Centerline{$U(x,\bover12\delta)\times U(\xi,\bover12\delta)
\subseteq U((x,\xi),\delta)\subseteq G$,}

\noindent sehingga $U(\xi,\bover12\delta)\in\Cal U_x$. Karena $\xi$
sebarang, $\Cal U_x$ merupakan selimut terbuka dari $[-n,n]$ dalam
$\Bbb R$. Menurut (i), selimut ini mempunyai subselimut berhingga, sebutlah
$U_0,\ldots,U_k$. Untuk setiap $j\le k$ kita dapat menemukan $H_j$, $G_j$
sedemikian sehingga $H_j$ merupakan subhimpunan terbuka dari $\BbbR^r$ yang
memuat $x$ dan $H_j\times U_j\subseteq G_j\in\Cal G$. Sekarang tetapkan
$H=\bigcap_{j\le k}H_j$. Ini merupakan subhimpunan terbuka dari $\BbbR^r$
yang memuat $x$, dan $H\times[-n,n]\subseteq\bigcup_{j\le k}G_j$ ditutupi
oleh suatu subkeluarga berhingga dari $\Cal G$. Jadi $x\in H\in\Cal H$.
Karena $x$ sebarang, $\Cal H$ menutupi $E$. \Qed

\medskip

\quad{\bf (iii)} Sekarang hipotesis induksi menyatakan bahwa $E$ kompak,
sehingga terdapat suatu subkeluarga berhingga $\Cal H_0$ dari $\Cal H$ yang
menutupi $E$. Untuk setiap $H\in\Cal H_0$, misalkan $\Cal G_H$ suatu
subkeluarga berhingga dari $\Cal G$ yang menutupi $H\times[-n,n]$. Maka
$\bigcup_{H\in\Cal H_0}\Cal G_H$ adalah subkeluarga berhingga dari $\Cal G$
yang menutupi $E\times[-n,n]=F$. Karena $\Cal G$ sebarang, $F$ kompak dan
induksi pun berlanjut.

\medskip

\quad{\bf (iv)} Jadi interval $[-\tbf{n},\tbf{n}]$ kompak dalam
$\BbbR^r$ untuk setiap $r$, $n$. Sekarang andaikan $K$ suatu himpunan
tertutup terbatas dalam $\BbbR^r$. Maka terdapat $n\in\Bbb N$ sedemikian
sehingga $K\subseteq[-\tbf{n},\tbf{n}]$, yaitu,
$K=K\cap[-\tbf{n},\tbf{n}]$. Karena $K$ tertutup dan
$[-\tbf{n},\tbf{n}]$ kompak, $K$ kompak menurut 2A2Ea.

Ini menyelesaikan bukti.
}%end of proof of 2A2F

\leader{2A2G}{Akibat} Jika $\phi:D\to\Bbb R$ kontinu, dengan
$D\subseteq\BbbR^r$, dan $K\subseteq D$ merupakan himpunan kompak tak kosong,
maka $\phi$ terbatas dan mencapai batas-batasnya pada $K$.

\proof{ Menurut 2A2Eb, $\phi[K]$ kompak; menurut 2A2F, himpunan itu tertutup
dan terbatas. Mengatakan bahwa $\phi[K]$ terbatas tepat berarti bahwa $\phi$
terbatas pada $K$. Karena $\phi[K]$ merupakan himpunan tak kosong dan terbatas,
himpunan itu mempunyai infimum $a$ dan supremum $b$; sekarang keduanya
termasuk dalam $\overline{\phi[K]}$ (menurut kriteria 2A2B(ii), atau dengan
cara lain); karena $\phi[K]$ tertutup, keduanya termasuk dalam $\phi[K]$,
yakni $\phi$ mencapai batas-batasnya.
}%end of proof of 2A2G

\leader{2A2H}{Meninjau kembali lim sup dan lim inf} Dalam \S1A3 saya
membahas secara singkat $\limsup_{n\to\infty}a_n$,
$\liminf_{n\to\infty}a_n$ untuk barisan real $\sequencen{a_n}$. Dalam
jilid ini kita memerlukan gagasan $\limsup_{\delta\downarrow 0}f(\delta)$,
$\liminf_{\delta\downarrow 0}f(\delta)$ untuk fungsi real $f$. Saya katakan
bahwa $\limsup_{\delta\downarrow 0}f(\delta)=u\in[-\infty,\infty]$ jika
(i) untuk setiap $v>u$ terdapat $\eta>0$ sedemikian sehingga $f(\delta)$
terdefinisi dan kurang dari atau sama dengan $v$ untuk setiap
$\delta\in\ocint{0,\eta}$ (ii) untuk setiap $v<u$ dan $\eta>0$ terdapat
$\delta\in\ocint{0,\eta}$ sedemikian sehingga $f(\delta)$ terdefinisi dan
lebih besar dari atau sama dengan $v$. Demikian pula,
$\liminf_{\delta\downarrow 0}f(\delta)=u\in[-\infty,\infty]$ jika (i) untuk
setiap $v<u$ terdapat $\eta>0$ sedemikian sehingga $f(\delta)$ terdefinisi
dan lebih besar dari atau sama dengan $v$ untuk setiap
$\delta\in\ocint{0,\eta}$ (ii) untuk setiap $v>u$ dan $\eta>0$ terdapat
$\delta\in\ocint{0,\eta}$ sedemikian sehingga $f(\delta)$ terdefinisi dan
kurang dari atau sama dengan $v$.

\leader{2A2I}{}\cmmnt{ Dalam kasus satu dimensi, kita mempunyai deskripsi
yang sangat sederhana bagi himpunan-himpunan terbuka.

\medskip

\noindent}{\bf Proposisi} Jika $G\subseteq\Bbb R$ sebarang himpunan terbuka,
maka himpunan itu dapat dinyatakan sebagai gabungan suatu keluarga terhitung dari
interval-interval terbuka yang saling lepas.

\proof{ Untuk $x$, $y\in G$, tuliskan $x\sim y$ jika $x\le y$ dan
$[x,y]\subseteq G$, atau $y\le x$ dan $[y,x]\subseteq G$. Mudah diperiksa
bahwa $\sim$ merupakan relasi ekuivalensi pada $G$. Misalkan $\Cal C$ adalah
himpunan kelas-kelas ekuivalensi terhadap $\sim$. Maka $\Cal C$ adalah suatu
partisi dari $G$. Sekarang setiap $C\in\Cal C$ adalah interval terbuka.
\Prf\ Tetapkan $a=\inf C$, $b=\sup C$ (dengan mengizinkan $a=-\infty$
dan/atau $b=\infty$ jika $C$ tak terbatas). Jika $a<x<b$, terdapat $y$,
$z\in C$ sedemikian sehingga $y\le x\le z$, sehingga
$[y,x]\subseteq[y,z]\subseteq G$ dan $y\sim x$ serta $x\in C$; jadi
$\ooint{a,b}\subseteq C$. Jika $x\in C$, terdapat interval terbuka $I$ yang
memuat $x$ dan termuat dalam $G$; karena $x\sim y$ untuk setiap $y\in I$,
$I\subseteq C$; jadi

\Centerline{$a\le\inf I<x<\sup I\le b$}

\noindent dan $x\in\ooint{a,b}$. Jadi $C=\ooint{a,b}$ adalah interval
terbuka.\ \Qed

Untuk melihat bahwa $\Cal C$ terhitung, perhatikan bahwa setiap anggota
$\Cal C$ memuat suatu anggota $\Bbb Q$, sehingga kita mempunyai suatu fungsi
surjektif dari subhimpunan $\Bbb Q$ ke $\Cal C$, dan $\Cal C$ terhitung
(1A1E).
}%end of proof of 2A2I

\discrpage
